\textbf{章节结构}：
\begin{itemize}
    \item \S 5.0 线性微分方程组解的存在唯一性定理
    \item \S 5.1 线性微分方程组的一般理论（齐次/非齐次）
    \item \S 5.2 常系数线性微分方程组（矩阵指数法）
    \item \S 5.3 拉普拉斯变换
\end{itemize}

\section{线性微分方程组解的存在唯一性定理}

\subsection{记号与定义}

\textbf{一阶微分方程组}：
$$\begin{cases} x'_1 = f_1(t, x_1, x_2, \ldots, x_n) \\ x'_2 = f_2(t, x_1, x_2, \ldots, x_n) \\ \vdots \\ x'_n = f_n(t, x_1, x_2, \ldots, x_n) \end{cases}$$
初值条件：$x_1(t_0) = \eta_1, x_2(t_0) = \eta_2, \ldots, x_n(t_0) = \eta_n$。

\textbf{一阶线性微分方程组}：
\begin{equation}x'_i = \sum_{j=1}^{n} a_{ij}(t)x_j + f_i(t), \quad i = 1, 2, \ldots, n \tag{5.1}\end{equation}
其中 $a_{ij}(t), f_i(t)$ 在 $[a,b]$ 上连续。

\textbf{矩阵向量形式}：令
$$\mathbf{A}(t) = \begin{pmatrix} a_{11}(t) & \cdots & a_{1n}(t) \\ \vdots & & \vdots \\ a_{n1}(t) & \cdots & a_{nn}(t) \end{pmatrix}, \quad \mathbf{f}(t) = \begin{pmatrix} f_1(t) \\ \vdots \\ f_n(t) \end{pmatrix}, \quad \mathbf{x} = \begin{pmatrix} x_1 \\ \vdots \\ x_n \end{pmatrix}$$
则
\begin{equation}\frac{d\mathbf{x}}{dt} = \mathbf{x}' = \mathbf{A}(t)\mathbf{x} + \mathbf{f}(t) \tag{5.4}\end{equation}

\textbf{矩阵与向量函数}：连续/可微/可积均按分量定义。运算性质：
\begin{itemize}
    \item $(\mathbf{A} + \mathbf{B})' = \mathbf{A}' + \mathbf{B}'$
    \item $(\mathbf{AB})' = \mathbf{A}'\mathbf{B} + \mathbf{A}\mathbf{B}'$
    \item $(\mathbf{Au})' = \mathbf{A}'\mathbf{u} + \mathbf{A}\mathbf{u}'$
\end{itemize}

\begin{mydefinition}{定义}
设 $\mathbf{A}(t)$ 是区间 $[\alpha, \beta]$ 上的连续 $n\times n$ 矩阵，$\mathbf{f}(t)$ 是 $n$ 维连续向量，则方程组 (5.4) 的解就是向量 $\mathbf{u}(t)$ 在区间 $[\alpha,\beta]$ 上连续且满足 $\mathbf{u}'(t) = \mathbf{A}(t)\mathbf{u}(t) + \mathbf{f}(t)$。
\end{mydefinition}

\begin{mydefinition}{初值问题}
Cauchy 问题
\begin{equation}\frac{d\mathbf{x}}{dt} = \mathbf{A}(t)\mathbf{x} + \mathbf{f}(t), \quad \mathbf{x}(t_0) = \boldsymbol{\eta} \tag{5.5}\end{equation}
的解就是方程组 (5.4) 在包含 $t_0$ 的区间 $[\alpha, \beta]$ 上的解 $\mathbf{u}(t)$，使得 $\mathbf{u}(t_0) = \boldsymbol{\eta}$。
\end{mydefinition}

\begin{example}
验证 $\mathbf{u}(t) = \begin{pmatrix} e^{-t} \\ -e^{-t} \end{pmatrix}$ 是初值问题 $\mathbf{x}' = \begin{pmatrix} 0 & 1 \\ 1 & 0 \end{pmatrix}\mathbf{x}, \ \mathbf{x}(0) = \begin{pmatrix} 1 \\ -1 \end{pmatrix}$ 在 $-\infty < t < +\infty$ 上的解。
\end{example}

\subsection{n 阶线性微分方程与一阶线性微分方程组等价}

\begin{example}
二阶方程 $x'' + p(t)x' + q(t)x = f(t)$ 令 $x_1 = x, x_2 = x'$，化为
$$\begin{cases} x'_1 = x_2 \\ x'_2 = -q(t)x_1 - p(t)x_2 + f(t) \end{cases} \iff \mathbf{x}' = \begin{pmatrix} 0 & 1 \\ -q(t) & -p(t) \end{pmatrix}\mathbf{x} + \begin{pmatrix} 0 \\ f(t) \end{pmatrix}$$
\end{example}

一般地，$n$ 阶线性方程 $x^{(n)} + a_1(t)x^{(n-1)} + \cdots + a_n(t)x = f(t)$ 令 $x_1 = x, x_2 = x', \ldots, x_n = x^{(n-1)}$，化为
\begin{equation}\mathbf{x}' = \begin{pmatrix} 0 & 1 & 0 & \cdots & 0 \\ 0 & 0 & 1 & \cdots & 0 \\ \vdots & & & \ddots & \vdots \\ 0 & 0 & 0 & \cdots & 1 \\ -a_n(t) & -a_{n-1}(t) & \cdots & \cdots & -a_1(t) \end{pmatrix}\mathbf{x} + \begin{pmatrix} 0 \\ \vdots \\ 0 \\ f(t) \end{pmatrix} \tag{5.7}\end{equation}

\begin{example}
将初值问题 $3x'' - 5x' + 2tx = \sin t, \ x(0) = 0, x'(0) = 0$ 化为一阶方程组初值问题。
\end{example}

\begin{solution}
令 $x_1 = x, x_2 = x'$，得 $\begin{cases} x'_1 = x_2 \\ x'_2 = \frac{5}{3}x_2 - \frac{2t}{3}x_1 + \frac{1}{3}\sin t \end{cases}$，$\mathbf{x}(0) = \begin{pmatrix} 0 \\ 0 \end{pmatrix}$。
\end{solution}

\begin{example}
将方程组 $\begin{cases} x'_1 = x_2 \\ x'_2 = -x_1 + x_2 \end{cases}$ 化为高阶方程。
\end{example}

\begin{solution}
由 $x_1 = x'_2 - x_2$ 消去，得 $x''_2 - x'_2 + x_2 = 0$。注意：\textbf{不是所有方程组都可化为高阶方程}。
\end{solution}

\subsection{存在唯一性定理}

\begin{myTheorem}{定理}
如果 $\mathbf{A}(t)$ 是 $n\times n$ 矩阵，$\mathbf{f}(t)$ 是 $n$ 维列向量，它们都在区间 $a \le t \le b$ 上连续，则对于区间上的任何数 $t_0$ 及任一常数向量 $\boldsymbol{\eta}$，方程组 (5.5) 存在唯一解 $\mathbf{x} = \boldsymbol{\varphi}(t)$，定义于整个区间 $a \le t \le b$ 上，且满足初始条件 $\mathbf{x}(t_0) = \boldsymbol{\eta}$。
\end{myTheorem}

\textbf{证明思路}（与第三章类似，逐次逼近法）：
\begin{enumerate}
    \item 初值问题等价于积分方程 $\mathbf{x}(t) = \boldsymbol{\eta} + \int_{t_0}^{t}[\mathbf{A}(s)\mathbf{x}(s) + \mathbf{f}(s)]ds$
    \item 构造 Picard 迭代向量函数序列 $\boldsymbol{\varphi}_0(t) = \boldsymbol{\eta}$，$\boldsymbol{\varphi}_{n+1}(t) = \boldsymbol{\eta} + \int_{t_0}^t[\mathbf{A}\boldsymbol{\varphi}_n + \mathbf{f}]ds$
    \item 证明序列在 $[a,b]$ 上一致收敛（Weierstrass 判别法，利用 $\|\mathbf{A}\|$ 有界）
    \item 极限是积分方程的连续解，即微分方程的解
    \item 唯一性：设另一解 $\boldsymbol{\psi}(t)$，令 $g(t) = \|\boldsymbol{\varphi}(t) - \boldsymbol{\psi}(t)\|$，用 Gronwall 型不等式证明 $g(t) \equiv 0$
\end{enumerate}

\begin{example}
求方程组初值问题 $\mathbf{x}' = \begin{pmatrix} -1 & 1 \\ 0 & 1 \end{pmatrix}\mathbf{x}, \ \mathbf{x}(0) = \begin{pmatrix} 1 \\ 0 \end{pmatrix}$ 的二次近似解。
\end{example}

\begin{solution}
Picard 迭代 $\boldsymbol{\varphi}_0(t) = \boldsymbol{\eta}$，$\boldsymbol{\varphi}_{n+1}(t) = \boldsymbol{\eta} + \int_{0}^{t}\mathbf{A}\boldsymbol{\varphi}_n(s)\,ds$：
$$\boldsymbol{\varphi}_1(t) = \begin{pmatrix} 1 - t \\ 0 \end{pmatrix}, \qquad \boldsymbol{\varphi}_2(t) = \begin{pmatrix} 1 - t + \frac{1}{2}t^2 \\ 0 \end{pmatrix}$$
\end{solution}

\subsection{简单方程组的消元法}

\begin{example}
求解方程组 $\begin{cases} x'_1 = -2x_1 + 3x_2 \\ x'_2 = x_1 - 2x_2 \end{cases}$。
\end{example}

\begin{solution}
保留 $x_2$，由第二式 $x_1 = x'_2 + 2x_2$，代入第一式消元得 $x''_2 + 4x'_2 + x_2 = 0$，特征方程 $\lambda^2 + 4\lambda + 1 = 0$ 的根为 $\lambda = -2 \pm \sqrt{3}$，解得
$$x_2 = c_1e^{(-2+\sqrt{3})t} + c_2e^{(-2-\sqrt{3})t}$$
代回 $x_1 = x'_2 + 2x_2$ 得
$$x_1 = \sqrt{3}\,c_1e^{(-2+\sqrt{3})t} - \sqrt{3}\,c_2e^{(-2-\sqrt{3})t}$$
即得方程组的通解。
\end{solution}

\begin{example}
求解方程组 $\begin{cases} \dfrac{dx}{dt} = y \\ \dfrac{dy}{dt} = \dfrac{x}{y} \end{cases}$
\end{example}

\begin{solution}
消去 $y$ 得 $x'' = \dfrac{x}{x'}$，令 $p = x'$ 降阶：$p' = \dfrac{x}{p}$，即 $p\dfrac{dp}{dx} = \dfrac{x}{p}$，故 $p^2\,dp = x\,dx$，积分得
$$p^3 = \frac{3}{2}x^2 + c_1$$
即 $\dfrac{dx}{dt} = \left(\frac{3}{2}x^2 + c_1\right)^{1/3}$，再分离变量积分即得隐式通解。
\end{solution}

\textbf{作业（5.0）}：p123：1（1、2）、2（1、2、3）、3

\section{线性微分方程组的一般理论}

一阶线性微分方程组：
\begin{equation}\frac{d\mathbf{x}}{dt} = \mathbf{A}(t)\mathbf{x} + \mathbf{f}(t) \tag{5.4}\end{equation}
其中 $\mathbf{A}(t), \mathbf{f}(t)$ 在 $a \le t \le b$ 上连续。若 $\mathbf{f}(t) \equiv 0$，则称
\begin{equation}\frac{d\mathbf{x}}{dt} = \mathbf{A}(t)\mathbf{x} \tag{5.8}\end{equation}
为一阶\textbf{齐线性}微分方程组；否则称 (5.4) 为\textbf{非齐线性}微分方程组。

\subsection{齐次线性微分方程组}

\begin{myTheorem}{叠加原理}
如果 $\mathbf{x}_1(t), \mathbf{x}_2(t), \ldots, \mathbf{x}_m(t)$ 是方程组 (5.8) 的 $m$ 个解，则它们的线性组合 $\mathbf{x}(t) = c_1\mathbf{x}_1(t) + \cdots + c_m\mathbf{x}_m(t)$ 也是 (5.8) 的解，其中 $c_1, \ldots, c_m$ 是任常数。
\end{myTheorem}

\textbf{定义（向量函数组线性相关/无关）}：设 $\mathbf{x}_1(t), \ldots, \mathbf{x}_m(t)$ 是定义在 $[a,b]$ 上的向量函数，若存在不全为零的常数 $c_1, \ldots, c_m$，使对所有 $a \le t \le b$ 有 $c_1\mathbf{x}_1(t) + \cdots + c_m\mathbf{x}_m(t) \equiv \mathbf{0}$，则称线性相关；否则线性无关。

\begin{example}
$\mathbf{x}_1 = \begin{pmatrix} \cos t \\ 1 \\ t \end{pmatrix}, \mathbf{x}_2 = \begin{pmatrix} 2\cos t \\ 2 \\ 2t \end{pmatrix}$ 在任何区间线性相关（取 $c_1 = 2, c_2 = -1$，则 $c_1\mathbf{x}_1 + c_2\mathbf{x}_2 \equiv \mathbf{0}$）。
\end{example}

\begin{example}
$\mathbf{x}_1 = \begin{pmatrix} e^t \\ e^{-t} \\ 0 \end{pmatrix}, \mathbf{x}_2 = \begin{pmatrix} e^{2t} \\ 0 \\ e^{3t} \end{pmatrix}, \mathbf{x}_3 = \begin{pmatrix} 0 \\ e^{3t} \\ e^{2t} \end{pmatrix}$ 在 $(-\infty, +\infty)$ 上线性无关（其 Wronsky 行列式 $W = -e^{7t} - e^{3t} \neq 0$）。
\end{example}

\subsubsection{Wronsky 行列式}

$n$ 个 $n$ 维向量函数构成的行列式
$$W[\mathbf{x}_1, \ldots, \mathbf{x}_n](t) = W(t) = \begin{vmatrix} x_{11}(t) & x_{12}(t) & \cdots & x_{1n}(t) \\ x_{21}(t) & x_{22}(t) & \cdots & x_{2n}(t) \\ \vdots & \vdots & & \vdots \\ x_{n1}(t) & x_{n2}(t) & \cdots & x_{nn}(t) \end{vmatrix}$$
称为这 $n$ 个向量函数构成的 \textbf{Wronsky 行列式}。

\begin{myTheorem}{定理}
如果向量函数 $\mathbf{x}_1(t), \ldots, \mathbf{x}_n(t)$ 在 $[a,b]$ 上线性相关，则它们的 Wronsky 行列式 $W(t) \equiv 0$（$a \le t \le b$）。
\end{myTheorem}

\begin{myTheorem}{定理}
如果 (5.8) 的解 $\mathbf{x}_1(t), \ldots, \mathbf{x}_n(t)$ 线性无关，则它们的 Wronsky 行列式 $W(t) \neq 0$（$a \le t \le b$）。
\end{myTheorem}

\begin{proof}
反证：若存在 $t_0$ 使 $W(t_0) = 0$，则数值向量组线性相关，存在不全为零常数使 $\sum c_i\mathbf{x}_i(t_0) = 0$；由叠加原理 $\mathbf{x}(t) = \sum c_i\mathbf{x}_i(t)$ 是 (5.8) 的解且 $\mathbf{x}(t_0) = 0$，由存在唯一性定理 $\mathbf{x}(t) \equiv 0$，与线性无关矛盾。
\end{proof}

\begin{remark}
\begin{itemize}
    \item \textbf{注 1}：(5.8) 的 $n$ 个解线性相关 $\iff W(t) \equiv 0$（$a \le t \le b$）
    \item \textbf{注 2}：(5.8) 的 $n$ 个解线性无关 $\iff W(t) \neq 0$（$a \le t \le b$）
\end{itemize}
即 (5.8) 的 $n$ 个解构成的 Wronsky 行列式\textbf{或者恒等于零，或者恒不等于零}。
\end{remark}

\begin{myTheorem}{定理}
(5.8) 一定存在 $n$ 个线性无关的解（取满足特殊单位初始条件的 $n$ 组解，$W(t_0) = 1 \neq 0$）。
\end{myTheorem}

\subsubsection{通解结构及基本解组}

\begin{myTheorem}{定理}
如果 $\mathbf{x}_1(t), \ldots, \mathbf{x}_n(t)$ 是 (5.8) 的 $n$ 个线性无关的解，则 (5.8) 的任一解均可表为
$$\mathbf{x}(t) = c_1\mathbf{x}_1(t) + \cdots + c_n\mathbf{x}_n(t)$$
其中 $c_1, \ldots, c_n$ 是相应的确定常数。
\end{myTheorem}

\textbf{推论 1}：(5.8) 的线性无关解的最大个数等于 $n$。

\textbf{基本解组}：(5.8) 的 $n$ 个线性无关解为 (5.8) 的一个基本解组。
\begin{itemize}
    \item 注 1：基本解组不唯一
    \item 注 2：(5.8) 所有解的集合构成一个 $n$ 维线性空间
    \item 注 3：由 n 阶方程初值问题与方程组初值问题的等价性，本章定理可平行推到 n 阶线性微分方程
\end{itemize}

\textbf{推论 2}：$n$ 阶微分方程 $x^{(n)} + a_1(t)x^{(n-1)} + \cdots + a_n(t)x = 0$（$a_i(t)$ 在 $[a,b]$ 上连续）若有 $n$ 个线性无关解，则任一解可表为其线性组合。

\subsubsection{解矩阵、基解矩阵及性质}

\begin{mydefinition}{定义}
如果一个 $n\times n$ 矩阵的每一列都是 (5.8) 的解，则称这个矩阵为 (5.8) 的\textbf{解矩阵}；如果该矩阵的列在 $[a,b]$ 上是 (5.8) 的线性无关解组，则称其为 (5.8) 的\textbf{基解矩阵}。以基本解组为列构成的矩阵 $\boldsymbol{\Phi}(t) = [\boldsymbol{\varphi}_1(t), \boldsymbol{\varphi}_2(t), \ldots, \boldsymbol{\varphi}_n(t)]$。
\end{mydefinition}

\begin{myTheorem}{定理}
(5.8) 一定存在一个基解矩阵；如果 $\boldsymbol{\Phi}(t)$ 是 (5.8) 的基解矩阵，那么 (5.8) 的任一解 $\boldsymbol{\psi}(t)$ 可表为
\begin{equation}\boldsymbol{\psi}(t) = \boldsymbol{\Phi}(t)\mathbf{C} \tag{5.15}\end{equation}
其中 $\mathbf{C}$ 是确定的 $n$ 维常向量。
\end{myTheorem}

\begin{myTheorem}{定理}
$\boldsymbol{\Phi}(t)$ 是 (5.8) 的基解矩阵的充要条件是：$\boldsymbol{\Phi}'(t) = \mathbf{A}(t)\boldsymbol{\Phi}(t)$，且存在 $t_0 \in [a,b]$ 使 $\det\boldsymbol{\Phi}(t_0) \neq 0$（此时 $\det\boldsymbol{\Phi}(t) \neq 0$ 对所有 $t$ 成立）。
\end{myTheorem}

\textbf{推论 1}：如果 $\boldsymbol{\Phi}(t)$ 是 (5.8) 在 $[a,b]$ 上的基解矩阵，$\mathbf{C}$ 是非奇异常数矩阵，那么 $\boldsymbol{\Phi}(t)\mathbf{C}$ 也是基解矩阵。

\textbf{推论 2}：如果 $\boldsymbol{\Phi}(t), \boldsymbol{\Psi}(t)$ 是 (5.8) 的两个基解矩阵，则存在非奇异常数矩阵 $\mathbf{C}$，使 $\boldsymbol{\Psi}(t) = \boldsymbol{\Phi}(t)\mathbf{C}$。

\begin{example}
验证 $\boldsymbol{\Phi}(t) = \begin{pmatrix} e^t & te^t \\ 0 & e^t \end{pmatrix}$ 是 $\mathbf{x}' = \begin{pmatrix} 1 & 1 \\ 0 & 1 \end{pmatrix}\mathbf{x}$ 的基解矩阵。
\end{example}

\begin{solution}
$\boldsymbol{\Phi}'(t) = \begin{pmatrix} e^t & (t+1)e^t \\ 0 & e^t \end{pmatrix} = \begin{pmatrix} 1 & 1 \\ 0 & 1 \end{pmatrix}\begin{pmatrix} e^t & te^t \\ 0 & e^t \end{pmatrix}$，且 $\det\boldsymbol{\Phi}(t) = e^{2t} \neq 0$，故是基解矩阵。
\end{solution}

\begin{example}
验证 $\boldsymbol{\Phi}(t) = \begin{pmatrix} e^t & e^{3t} \\ -e^t & e^{3t} \end{pmatrix}$ 是 $\mathbf{x}' = \begin{pmatrix} 2 & 1 \\ 1 & 2 \end{pmatrix}\mathbf{x}$ 的基解矩阵，并求其通解。
\end{example}

\begin{solution}
$\det\boldsymbol{\Phi}(t) = 2e^{4t} \neq 0$，且每列 $e^t\begin{pmatrix}1\\-1\end{pmatrix}, e^{3t}\begin{pmatrix}1\\1\end{pmatrix}$ 代入可验证是解。通解 $\mathbf{x}(t) = \boldsymbol{\Phi}(t)\mathbf{C} = c_1\begin{pmatrix}e^t\\-e^t\end{pmatrix} + c_2\begin{pmatrix}e^{3t}\\e^{3t}\end{pmatrix}$。
\end{solution}

\textbf{作业（5.1.1）}：p123：4、5、6、7、8

\subsection{非齐次线性微分方程组}

\textbf{性质 1}：如果 $\boldsymbol{\varphi}(t)$ 是 (5.4) 的解，$\boldsymbol{\psi}(t)$ 是对应齐次方程组 (5.8) 的解，则 $\boldsymbol{\varphi}(t) + \boldsymbol{\psi}(t)$ 是 (5.4) 的解。

\textbf{性质 2}：如果 $\boldsymbol{\varphi}(t), \bar{\boldsymbol{\varphi}}(t)$ 是 (5.4) 的两个解，则 $\boldsymbol{\varphi}(t) - \bar{\boldsymbol{\varphi}}(t)$ 是 (5.8) 的解。

\textbf{性质 3}：设 $\mathbf{f}(t) = \mathbf{f}_1(t) + \cdots + \mathbf{f}_m(t)$，且 $\mathbf{x}_j$ 是 $\mathbf{x}' = \mathbf{A}(t)\mathbf{x} + \mathbf{f}_j(t)$ 的解，则 $\mathbf{x} = \sum \mathbf{x}_j$ 是 (5.4) 的解。

\begin{myTheorem}{通解结构}
设 $\boldsymbol{\Phi}(t)$ 是 (5.8) 的基解矩阵，$\bar{\boldsymbol{\varphi}}(t)$ 是 (5.4) 的某一解，则 (5.4) 的任一解可表为
$$\boldsymbol{\varphi}(t) = \boldsymbol{\Phi}(t)\mathbf{C} + \bar{\boldsymbol{\varphi}}(t)$$
其中 $\mathbf{C}$ 是确定的常数列向量。
\end{myTheorem}

\subsubsection{常数变易公式}

设 $\boldsymbol{\Phi}(t)$ 是 (5.8) 的基解矩阵，则 (5.8) 的通解为 $\mathbf{x}(t) = \boldsymbol{\Phi}(t)\mathbf{C}$。寻求 (5.4) 形如
\begin{equation}\mathbf{x}(t) = \boldsymbol{\Phi}(t)\mathbf{C}(t) \tag{5.17}\end{equation}
的解，代入 (5.4) 得 $\mathbf{C}'(t) = \boldsymbol{\Phi}^{-1}(t)\mathbf{f}(t)$，从 $t_0$ 积分并取 $\mathbf{C}(t_0) = 0$，得

\begin{myTheorem}{定理}
如果 $\boldsymbol{\Phi}(t)$ 是 (5.8) 的基解矩阵，则：
\begin{enumerate}
    \item 向量函数 $\boldsymbol{\varphi}(t) = \boldsymbol{\Phi}(t)\int_{t_0}^{t}\boldsymbol{\Phi}^{-1}(s)\mathbf{f}(s)\,ds$ 是 (5.4) 的解，且满足初始条件 $\boldsymbol{\varphi}(t_0) = 0$
    \item 方程组 (5.4) 的通解为 $\mathbf{x}(t) = \boldsymbol{\Phi}(t)\mathbf{C} + \boldsymbol{\Phi}(t)\int_{t_0}^{t}\boldsymbol{\Phi}^{-1}(s)\mathbf{f}(s)\,ds$
\end{enumerate}
\end{myTheorem}

\begin{remark}
\textbf{注 1}：满足初始条件 $\boldsymbol{\varphi}(t_0) = \boldsymbol{\eta}$ 的解为
\begin{equation}\boldsymbol{\varphi}(t) = \boldsymbol{\Phi}(t)\boldsymbol{\Phi}^{-1}(t_0)\boldsymbol{\eta} + \boldsymbol{\Phi}(t)\int_{t_0}^{t}\boldsymbol{\Phi}^{-1}(s)\mathbf{f}(s)\,ds \tag{5.20}\end{equation}
\textbf{注 2}：公式 (5.19) 或 (5.20) 称为 (5.4) 的\textbf{常数变易公式}。
\end{remark}

\begin{example}
求初值问题 $\mathbf{x}' = \begin{pmatrix} 1 & 1 \\ 0 & 1 \end{pmatrix}\mathbf{x} + \begin{pmatrix} e^{-t} \\ 0 \end{pmatrix}, \ \mathbf{x}(0) = \begin{pmatrix} -1 \\ 1 \end{pmatrix}$ 的解。
\end{example}

\begin{solution}
由例 3 基解矩阵 $\boldsymbol{\Phi}(t) = \begin{pmatrix} e^t & te^t \\ 0 & e^t \end{pmatrix}$，$\boldsymbol{\Phi}^{-1}(s) = \begin{pmatrix} e^{-s} & -se^{-s} \\ 0 & e^{-s} \end{pmatrix}$，代入常数变易公式计算得
$$\boldsymbol{\varphi}(t) = \begin{pmatrix} te^t - e^t + \frac{1}{2}(e^t - e^{-t}) \\ e^t \end{pmatrix}$$
\end{solution}

\subsubsection{n 阶线性微分方程的常数变易公式}

设 $x_1(t), \ldots, x_n(t)$ 是 (5.21) 对应齐次方程的基本解组，则非齐方程 (5.21) 满足零初始条件 $x(t_0) = x'(t_0) = \cdots = x^{(n-1)}(t_0) = 0$ 的解为
\begin{equation}\varphi(t) = \sum_{k=1}^{n} x_k(t)\int_{t_0}^{t} \frac{W[x_1(s), \ldots, x_{k-1}(s), \mathbf{0}, x_{k+1}(s), \ldots, x_n(s)]}{W[x_1(s), \ldots, x_n(s)]} f(s)\,ds \tag{5.22}\end{equation}
公式 (5.22) 称为 (5.21) 的\textbf{常数变易公式}。方程 (5.21) 的通解为
$$x(t) = c_1x_1(t) + c_2x_2(t) + \cdots + c_nx_n(t) + \varphi(t)$$

\textbf{$n = 2$ 时}的常数变易公式：
\begin{equation}\varphi(t) = \int_{t_0}^{t} \frac{x_1(s)x_2(t) - x_2(s)x_1(t)}{W[x_1(s), x_2(s)]}f(s)\,ds \tag{5.24}\end{equation}
通解 $x(t) = c_1x_1(t) + c_2x_2(t) + \varphi(t)$。

\begin{example}
求方程 $x'' + x = \tan t$ 的一个解。
\end{example}

\begin{solution}
基本解组 $x_1 = \cos t, x_2 = \sin t$，$W = \cos^2 t + \sin^2 t = 1$。由 (5.24)：
$$\varphi(t) = \int_0^t (\cos s\sin t - \sin s\cos t)\tan s\,ds = \sin t\int_0^t \sin s\,ds - \cos t\int_0^t \sin s\tan s\,ds$$
$$= \sin t(1 - \cos t) - \cos t(\sin t - \ln|\sec t + \tan t|)$$
注意到 $-\sin t$ 是对应齐次方程的解，故 $\varphi(t) = -\cos t\ln|\sec t + \tan t|$ 也是原方程的一个解。
\end{solution}

\textbf{作业（5.1.2）}：p124-125：9、10、11、12、13、14、16、17

\section{常系数线性微分方程组}

一阶常系数线性微分方程组：$\mathbf{x}' = \mathbf{A}\mathbf{x} + \mathbf{f}(t)$，系数矩阵 $\mathbf{A}$ 为 $n\times n$ 常数矩阵；若 $\mathbf{f}(t) = 0$，对应齐次方程组为
\begin{equation}\frac{d\mathbf{x}}{dt} = \mathbf{A}\mathbf{x} \tag{5.26}\end{equation}
本节主要讨论 (5.26) 的基解矩阵的求法。

\subsection{矩阵指数 exp A 的定义和求法}

\begin{mydefinition}{定义}
设 $\mathbf{A}$ 为 $n\times n$ 常数矩阵，定义矩阵指数为矩阵级数的和
\begin{equation}\exp\mathbf{A} = \sum_{k=0}^{\infty}\frac{\mathbf{A}^k}{k!} = \mathbf{E} + \mathbf{A} + \frac{\mathbf{A}^2}{2!} + \cdots + \frac{\mathbf{A}^m}{m!} + \cdots \tag{5.27}\end{equation}
其中 $\mathbf{E}$ 为单位矩阵，$\mathbf{A}^0 = \mathbf{E}$。
\end{mydefinition}

\begin{itemize}
    \item \textbf{注 1}：矩阵级数收敛（$\left\|\frac{\mathbf{A}^k}{k!}\right\| \le \frac{\|\mathbf{A}\|^k}{k!}$，数项级数 $\sum \frac{\|\mathbf{A}\|^k}{k!}$ 收敛）
    \item \textbf{注 2}：级数 $\exp(\mathbf{A}t) = \sum_{k=0}^{\infty}\frac{\mathbf{A}^k t^k}{k!}$ 在 $t$ 的任何有限区间上一致收敛
\end{itemize}

\textbf{矩阵指数的性质}：
\begin{enumerate}
    \item 若 $\mathbf{AB} = \mathbf{BA}$，则 $e^{\mathbf{A}+\mathbf{B}} = e^{\mathbf{A}}e^{\mathbf{B}}$
    \item 对任何矩阵 $\mathbf{A}$，$(\exp\mathbf{A})^{-1}$ 存在，且 $(\exp\mathbf{A})^{-1} = \exp(-\mathbf{A})$
    \item 若 $\mathbf{T}$ 非奇异，则 $\exp(\mathbf{T}^{-1}\mathbf{A}\mathbf{T}) = \mathbf{T}^{-1}(\exp\mathbf{A})\mathbf{T}$
\end{enumerate}

\subsubsection{常系数齐线性微分方程组的基解矩阵}

\begin{myTheorem}{定理}
矩阵 $\boldsymbol{\Phi}(t) = \exp(\mathbf{A}t)$ 是 (5.26) 的基解矩阵，且 $\boldsymbol{\Phi}(0) = \mathbf{E}$。
\end{myTheorem}

\begin{proof}
$\boldsymbol{\Phi}'(t) = \mathbf{A}\exp(\mathbf{A}t) = \mathbf{A}\boldsymbol{\Phi}(t)$，且 $\boldsymbol{\Phi}(0) = \mathbf{E}$，$\det\boldsymbol{\Phi}(0) = 1 \neq 0$。
\end{proof}

\begin{example}[对角矩阵]
$\mathbf{A} = \mathrm{diag}(a_1, a_2, \ldots, a_n)$，基解矩阵 $\boldsymbol{\Phi}(t) = \exp(\mathbf{A}t) = \mathrm{diag}(e^{a_1t}, e^{a_2t}, \ldots, e^{a_nt})$。
\end{example}

\begin{example}[Jordan 块矩阵]
$\mathbf{A} = \begin{pmatrix} 2 & 1 \\ 0 & 2 \end{pmatrix} = \mathbf{D} + \mathbf{N}$，其中 $\mathbf{D} = 2\mathbf{E}$，$\mathbf{N} = \begin{pmatrix} 0 & 1 \\ 0 & 0 \end{pmatrix}$ 幂零（$\mathbf{N}^2 = 0$），$\mathbf{DN} = \mathbf{ND}$。则
$$\exp(\mathbf{A}t) = \exp(\mathbf{D}t)\exp(\mathbf{N}t) = e^{2t}\mathbf{E}\cdot(\mathbf{E} + \mathbf{N}t) = \begin{pmatrix} e^{2t} & te^{2t} \\ 0 & e^{2t} \end{pmatrix}$$
\end{example}

\textbf{基解矩阵的一种求法}：设 $\mathbf{A} = \mathbf{T}^{-1}\mathbf{J}\mathbf{T}$（$\mathbf{J}$ 为 Jordan 矩阵），则 $e^{\mathbf{A}t} = \mathbf{T}^{-1}e^{\mathbf{J}t}\mathbf{T}$，其中 $\mathbf{J} = \mathrm{diag}(\mathbf{J}_1, \ldots, \mathbf{J}_k)$，$e^{\mathbf{J}t} = \mathrm{diag}(e^{\mathbf{J}_1t}, \ldots, e^{\mathbf{J}_kt})$。注：$e^{\mathbf{A}t}\mathbf{T}^{-1} = \mathbf{T}^{-1}e^{\mathbf{J}t}$ 也是基解矩阵。

\subsubsection{基解矩阵的计算公式}

\textbf{特征值与特征向量的关系}：寻求 (5.26) 形如 $\boldsymbol{\varphi}(t) = e^{\lambda t}\mathbf{c}$（$\mathbf{c} \neq 0$）的解，代入得 $(\lambda\mathbf{E} - \mathbf{A})\mathbf{c} = 0$，有非零解 $\iff \det(\lambda\mathbf{E} - \mathbf{A}) = 0$。

\textbf{结论}：(5.26) 有非零解 $\boldsymbol{\varphi}(t) = e^{\lambda t}\mathbf{c}$ 的充要条件是 $\lambda$ 是 $\mathbf{A}$ 的特征根，$\mathbf{c}$ 是对应的特征向量。

\begin{example}
求 $\mathbf{A} = \begin{pmatrix} 3 & 5 \\ -5 & 3 \end{pmatrix}$ 的特征值和特征向量。特征方程 $\lambda^2 - 6\lambda + 34 = 0$，$\lambda_{1,2} = 3 \pm 5i$；$\lambda_1$ 对应特征向量 $u = \alpha\begin{pmatrix}1\\i\end{pmatrix}$，$\lambda_2$ 对应 $v = \beta\begin{pmatrix}i\\1\end{pmatrix}$。
\end{example}

\begin{example}
$\mathbf{A} = \begin{pmatrix} 2 & 1 \\ -1 & 4 \end{pmatrix}$，特征值 $\lambda = 3$（二重根），特征向量 $c = \alpha\begin{pmatrix}1\\1\end{pmatrix}$。
\end{example}

\subsubsection{(1) 矩阵 A 具有 n 个线性无关的特征向量时}

\begin{myTheorem}{定理}
如果矩阵 $\mathbf{A}$ 具有 $n$ 个线性无关的特征向量 $\mathbf{v}_1, \mathbf{v}_2, \ldots, \mathbf{v}_n$，相应特征值为 $\lambda_1, \lambda_2, \ldots, \lambda_n$（不必互不相同），则矩阵
$$\boldsymbol{\Phi}(t) = [e^{\lambda_1 t}\mathbf{v}_1, e^{\lambda_2 t}\mathbf{v}_2, \ldots, e^{\lambda_n t}\mathbf{v}_n], \quad -\infty < t < +\infty$$
是 (5.26) 的一个基解矩阵。
\end{myTheorem}

\begin{proof}
每个 $e^{\lambda_j t}\mathbf{v}_j$ 都是解；$\det\boldsymbol{\Phi}(0) = \det[\mathbf{v}_1, \ldots, \mathbf{v}_n] \neq 0$。
\end{proof}

\begin{example}
$\mathbf{x}' = \begin{pmatrix} 3 & 5 \\ -5 & 3 \end{pmatrix}\mathbf{x}$ 的基解矩阵为 $\boldsymbol{\Phi}(t) = \begin{pmatrix} e^{(3+5i)t} & ie^{(3-5i)t} \\ ie^{(3+5i)t} & e^{(3-5i)t} \end{pmatrix}$。
\end{example}

\textbf{注与例 6（实基解矩阵）}：一般来说 $\boldsymbol{\Phi}(t)$ 不一定是 $\exp(\mathbf{A}t)$；但 $\exp(\mathbf{A}t) = \boldsymbol{\Phi}(t)\boldsymbol{\Phi}^{-1}(0)$。例 5 的实基解矩阵：
$$\exp(\mathbf{A}t) = e^{3t}\begin{pmatrix} \cos 5t & \sin 5t \\ -\sin 5t & \cos 5t \end{pmatrix}$$
满足 $\varphi(0) = \begin{pmatrix}1\\1\end{pmatrix}$ 的解为 $\varphi(t) = e^{3t}\begin{pmatrix}\cos 5t + \sin 5t \\ \cos 5t - \sin 5t\end{pmatrix}$。

\begin{example}
$\mathbf{x}' = \begin{pmatrix} 5 & -28 & -18 \\ -1 & 5 & 3 \\ 3 & -16 & -10 \end{pmatrix}\mathbf{x}$，特征方程 $\lambda(1-\lambda^2) = 0$，根 $\lambda = 0, 1, -1$，对应特征向量 $v_1 = \begin{pmatrix}-2\\-1\\1\end{pmatrix}, v_2 = \begin{pmatrix}2\\-1\\2\end{pmatrix}, v_3 = \begin{pmatrix}3\\0\\1\end{pmatrix}$，通解 $\varphi(t) = c_1v_1 + c_2v_2e^t + c_3v_3e^{-t}$。
\end{example}

\subsubsection{(2) 矩阵 A 的特征根有重根时}

假设 $\mathbf{A}$ 的特征值为 $\lambda_1, \ldots, \lambda_k$，重数分别为 $n_1, \ldots, n_k$，$n_1 + \cdots + n_k = n$。由高等代数，$n$ 维空间 $\mathbf{U}$ 的子集
\begin{equation}\mathbf{U}_j = \{\mathbf{u} \in \mathbf{U} \mid (\mathbf{A} - \lambda_j\mathbf{E})^{n_j}\mathbf{u} = 0\} \tag{5.36}\end{equation}
是 $\mathbf{U}$ 的 $n_j$ 维不变子空间，且 $\mathbf{U} = \mathbf{U}_1 \oplus \mathbf{U}_2 \oplus \cdots \oplus \mathbf{U}_k$。将 $\boldsymbol{\eta}$ 分解为 $\boldsymbol{\eta} = \mathbf{v}_1 + \cdots + \mathbf{v}_k$（$\mathbf{v}_j \in \mathbf{U}_j$），有 $(\mathbf{A} - \lambda_j\mathbf{E})^l\mathbf{v}_j = 0$（$l \ge n_j$）。

\textbf{解的表达式}：满足初始条件 $\varphi(0) = \boldsymbol{\eta}$ 的解可写成\textbf{有限项组合}：
\begin{equation}\varphi(t) = (\exp\mathbf{A}t)\boldsymbol{\eta} = \sum_{j=1}^{k} e^{\lambda_j t}\left\{\sum_{i=0}^{n_j - 1}\frac{(\mathbf{A} - \lambda_j\mathbf{E})^i t^i}{i!}\right\}\mathbf{v}_j \tag{5.40}\end{equation}

\begin{remark}
\textbf{注 1}：当 $\mathbf{A}$ 只有一个特征值时，对任何 $\mathbf{u}$ 有 $(\mathbf{A} - \lambda\mathbf{E})^n\mathbf{u} = 0$，故
\begin{equation}\exp(\mathbf{A}t) = e^{\lambda t}\sum_{i=0}^{n-1}\frac{(\mathbf{A} - \lambda\mathbf{E})^i t^i}{i!} \tag{5.41}\end{equation}
\textbf{注 2}：为从 (5.40) 求 $\exp(\mathbf{A}t)$，令 $\boldsymbol{\eta} = \mathbf{e}_1, \ldots, \mathbf{e}_n$（单位向量），求得 $n$ 个线性无关的解，以这些解为列即得 $\exp(\mathbf{A}t)$。
\end{remark}

\begin{example}
$\mathbf{x}' = \begin{pmatrix} 2 & 1 \\ -1 & 4 \end{pmatrix}\mathbf{x}, \ \varphi(0) = \boldsymbol{\eta}$，并求 $\exp(\mathbf{A}t)$。
\end{example}

\begin{solution}
$\lambda = 3$ 为二重特征值，只有一个子空间 $\mathbf{U}_1$。由 (5.40)：
$$\varphi(t) = e^{3t}\{\mathbf{E} + t(\mathbf{A} - 3\mathbf{E})\}\boldsymbol{\eta} = e^{3t}\begin{pmatrix} 1-t & t \\ -t & 1+t \end{pmatrix}\boldsymbol{\eta}$$
故 $\exp(\mathbf{A}t) = e^{3t}\begin{pmatrix} 1-t & t \\ -t & 1+t \end{pmatrix}$。
\end{solution}

\begin{example}
$\mathbf{A}$ 为 5 阶矩阵（$\lambda = -4$ 五重特征值），$(A+4E)^3 = 0$，由 (5.41)：
$$\exp(\mathbf{A}t) = e^{-4t}\left\{\mathbf{E} + t(\mathbf{A}+4\mathbf{E}) + \frac{t^2}{2!}(\mathbf{A}+4\mathbf{E})^2\right\} = e^{-4t}\begin{pmatrix} 1 & t & t^2/2 & 0 & 0 \\ 0 & 1 & t & 0 & 0 \\ 0 & 0 & 1 & 0 & 0 \\ 0 & 0 & 0 & 1 & 0 \\ 0 & 0 & 0 & 0 & 1 \end{pmatrix}$$
\end{example}

\begin{example}
$\begin{cases} x'_1 = 3x_1 - x_2 + x_3 \\ x'_2 = 2x_1 + x_3 \\ x'_3 = x_1 - x_2 + 2x_3 \end{cases}$，$\mathbf{A} = \begin{pmatrix} 3 & -1 & 1 \\ 2 & 0 & 1 \\ 1 & -1 & 2 \end{pmatrix}$，特征方程 $(\lambda-1)(\lambda-2)^2 = 0$。
\begin{itemize}
    \item $\mathbf{U}_1$（$\lambda=1$）：$(\mathbf{A}-\mathbf{E})\mathbf{u} = 0$ 解为 $\mathbf{u}_1 = \alpha\begin{pmatrix}0\\1\\1\end{pmatrix}$
    \item $\mathbf{U}_2$（$\lambda=2$）：$(\mathbf{A}-2\mathbf{E})^2\mathbf{u} = 0$ 解为 $\mathbf{u}_2 = \begin{pmatrix}\beta\\\beta\\\gamma\end{pmatrix}$
\end{itemize}
分解 $\boldsymbol{\eta} = \mathbf{v}_1 + \mathbf{v}_2$ 得 $\mathbf{v}_1 = \begin{pmatrix}0\\\eta_2-\eta_1\\\eta_2-\eta_1\end{pmatrix}$，$\mathbf{v}_2 = \begin{pmatrix}\eta_1\\\eta_1\\\eta_3-\eta_2+\eta_1\end{pmatrix}$。解为
$$\varphi(t) = e^t\mathbf{v}_1 + e^{2t}\{\mathbf{E} + t(\mathbf{A}-2\mathbf{E})\}\mathbf{v}_2$$
令 $\boldsymbol{\eta}$ 取三个单位向量得 $\exp(\mathbf{A}t)$（见讲义 p37 的 $3\times3$ 矩阵结果）。
\end{example}

\subsubsection{(3) 非齐线性方程的解}

非齐方程组 $\mathbf{x}' = \mathbf{A}\mathbf{x} + \mathbf{f}(t)$ (5.43) 满足初始条件 $\varphi(t_0) = \boldsymbol{\eta}$ 的解：
$$\varphi(t) = \exp[\mathbf{A}(t-t_0)]\boldsymbol{\eta} + \int_{t_0}^{t}\exp[\mathbf{A}(t-s)]\mathbf{f}(s)\,ds$$
或写为基解矩阵形式 $\varphi(t) = \boldsymbol{\Phi}(t)\boldsymbol{\Phi}^{-1}(t_0)\boldsymbol{\eta} + \boldsymbol{\Phi}(t)\int_{t_0}^{t}\boldsymbol{\Phi}^{-1}(s)\mathbf{f}(s)\,ds$。

\begin{example}
$\mathbf{A} = \begin{pmatrix} 3 & 5 \\ -5 & 3 \end{pmatrix}, \mathbf{f}(t) = \begin{pmatrix} e^{-t} \\ 0 \end{pmatrix}$，$\varphi(0) = \begin{pmatrix}0\\1\end{pmatrix}$。
\end{example}

\begin{solution}
$\exp(\mathbf{A}t) = e^{3t}\begin{pmatrix}\cos 5t & \sin 5t \\ -\sin 5t & \cos 5t\end{pmatrix}$，代入得
$$\varphi(t) = \frac{1}{41}e^{3t}\begin{pmatrix} 4\cos 5t + 46\sin 5t - 4e^{-4t} \\ 46\cos 5t - 4\sin 5t - 5e^{-4t} \end{pmatrix}$$
\end{solution}

\textbf{作业（5.2）}：p137-138：1（1、2）、2、3（2、3）、4（1、4）、5（3）、6（2、3）、7、8（1、2、3）、9

\section{拉普拉斯变换}

\subsection{拉普拉斯变换}

\textbf{定义}：
$$\mathcal{L}[f(t)] \equiv F(s) = \int_0^{\infty} e^{-st}f(t)\,dt, \quad \mathrm{Re}\,s > \delta$$
其中 $f(t)$ 在 $t \ge 0$ 有定义，且满足 $|f(t)| \le Me^{\delta t}$（$M, \delta$ 为正数）。称 $f(t)$ 为\textbf{原函数}，$F(s)$ 为\textbf{像函数}。

\begin{example}
$\mathcal{L}[t^n] = \dfrac{n!}{s^{n+1}}$（递推：$\mathcal{L}[t^n] = \dfrac{n}{s}\mathcal{L}[t^{n-1}]$）。
\end{example}

\begin{example}
$\mathcal{L}[e^{at}] = \dfrac{1}{s-a}$（$s > a$）。
\end{example}

\textbf{常用变换对}：$\mathcal{L}[\cos\omega t] = \dfrac{s}{s^2+\omega^2}$，$\mathcal{L}[\sin\omega t] = \dfrac{\omega}{s^2+\omega^2}$。

\subsubsection{拉普拉斯变换的性质}

\begin{enumerate}
    \item \textbf{$n$ 阶导数的拉普拉斯变换}：
    $$\mathcal{L}[f^{(n)}(t)] = s^n\mathcal{L}[f(t)] - s^{n-1}f(0) - s^{n-2}f'(0) - \cdots - f^{(n-1)}(0)$$
    \item \textbf{其他性质}（若 $\mathcal{L}[f(t)] \equiv F(s)$）：
    \begin{itemize}
        \item $\dfrac{d}{ds}F(s) = -\mathcal{L}[tf(t)]$，$\dfrac{d^n}{ds^n}F(s) = (-1)^n\mathcal{L}[t^nf(t)]$
        \item $\mathcal{L}[e^{at}f(t)] = F(s-a)$
    \end{itemize}
\end{enumerate}

\subsubsection{应用：求解微分方程}

给定 $x^{(n)} + a_1x^{(n-1)} + \cdots + a_nx = f(t)$，初始条件 $x(0) = x_0, x'(0) = x_0^{(1)}, \ldots$。对方程两端作拉普拉斯变换，设 $X(s) = \mathcal{L}[x(t)], F(s) = \mathcal{L}[f(t)]$，得
$$A(s)X(s) = F(s) + B(s)$$
其中 $A(s) = s^n + a_1s^{n-1} + \cdots + a_n$ 为特征多项式，$B(s)$ 由初始条件决定。则 $X(s) = \dfrac{F(s) + B(s)}{A(s)}$，解为 $x(t) = \mathcal{L}^{-1}[X(s)]$。

\begin{example}
求解 $x'' + 2x' + x = e^{-t}, \ x(1) = x'(1) = 0$。
\end{example}

\begin{solution}
令 $\tau = t - 1$，$u(\tau) = x(t)$，得 $u'' + 2u' + u = e^{-(\tau+1)}, \ u(0) = u'(0) = 0$。变换得 $(s^2+2s+1)U(s) = \dfrac{1}{e}\cdot\dfrac{1}{s+1}$，$U(s) = \dfrac{1}{e(s+1)^3}$，反变换 $u(\tau) = \dfrac{1}{2e}\tau^2e^{-\tau}$，故 $x(t) = \dfrac{1}{2}(t-1)^2e^{-t}$。
\end{solution}

\begin{example}
求解 $x'' - x = 4\cos t + 5\sin 2t, \ x(0) = -1, x'(0) = -2$。
\end{example}

\begin{solution}
变换得 $s^2X(s) + s + 2 - X(s) = \dfrac{4}{s^2+1} + \dfrac{10}{s^2+4}$，解得 $X(s) = -\dfrac{2}{s^2+1} - \dfrac{s}{s^2+4}$，反变换得 $x(t) = -2\sin t - \cos 2t$。
\end{solution}

\subsection{微分方程组的拉普拉斯变换}

\textbf{定义}：$\mathcal{L}[\mathbf{f}(t)] = \int_0^{\infty} e^{-st}\mathbf{f}(t)\,dt$（按分量定义）。

\begin{myTheorem}{定理}
若向量函数 $\mathbf{f}(t)$ 满足 $\|\mathbf{f}(t)\| \le Me^{\delta t}$（对所有充分大的 $t$），则初值问题 $\mathbf{x}' = \mathbf{A}\mathbf{x} + \mathbf{f}(t), \ \mathbf{x}(0) = \boldsymbol{\eta}$ 的解 $\varphi(t)$ 及其导数 $\varphi'(t)$ 均满足类似不等式，从而其拉普拉斯变换都存在。
\end{myTheorem}

\textbf{推论}：若 $|f(t)| \le Me^{\delta t}$，则 $x^{(n)} + a_1x^{(n-1)} + \cdots + a_nx = f(t)$ 初值问题的解及其直到 $n$ 阶导数均存在拉普拉斯变换。

\begin{example}
用拉普拉斯变换求解 $\begin{cases} x'_1 = 3x_1 + 5x_2 + e^{-t} \\ x'_2 = -5x_1 + 3x_2 \end{cases}$，$\varphi_1(0) = 0, \varphi_2(0) = 1$。
\end{example}

\begin{solution}
令 $X_1 = \mathcal{L}[\varphi_1], X_2 = \mathcal{L}[\varphi_2]$，变换得
$$\begin{cases} (s-3)X_1 - 5X_2 = \dfrac{1}{s+1} \\ 5X_1 + (s-3)X_2 = 1 \end{cases}$$
解得 $X_1(s) = \dfrac{1}{41}\left[\dfrac{4(s-3)}{(s-3)^2+25} + \dfrac{46\cdot 5}{(s-3)^2+25} - \dfrac{4}{s+1}\right]$，反变换得
$$\varphi_1(t) = \frac{1}{41}e^{3t}(4\cos 5t + 46\sin 5t - 4e^{-4t}), \quad \varphi_2(t) = \frac{1}{41}e^{3t}(46\cos 5t - 4\sin 5t - 5e^{-4t})$$
\end{solution}

\begin{example}
$\begin{cases} x'_1 = 2x_1 + x_2 \\ x'_2 = -x_1 + 4x_2 \end{cases}$，$\varphi_1(0) = 0, \varphi_2(0) = 1$。
\end{example}

\begin{solution}
变换得 $\begin{cases} (s-2)X_1 - X_2 = 0 \\ X_1 + (s-4)X_2 = 1 \end{cases}$，解得 $X_1 = \dfrac{1}{(s-3)^2}, X_2 = \dfrac{1}{s-3} + \dfrac{1}{(s-3)^2}$，反变换
$$\varphi_1(t) = te^{3t}, \quad \varphi_2(t) = (1+t)e^{3t}$$
\end{solution}

\begin{example}
$\begin{cases} x''_1 - 2x'_1 - x'_2 + 2x_2 = 0 \\ x'_1 - 2x_1 + x'_2 = -2e^{-t} \end{cases}$，$\varphi_1(0) = 3, \varphi'_1(0) = 2, \varphi_2(0) = 0$。
\end{example}

\begin{solution}
变换得 $\begin{cases} (s^2-2s)X_1 - (s-2)X_2 = 3s - 4 \\ (s-2)X_1 + sX_2 = \dfrac{3s+1}{s+1} \end{cases}$，解得 $X_1 = \dfrac{1}{s-1} + \dfrac{1}{s+1} + \dfrac{1}{s-2}$，$X_2 = \dfrac{1}{s-1} - \dfrac{1}{s+1}$，反变换
$$\varphi_1(t) = e^t + e^{-t} + e^{2t}, \quad \varphi_2(t) = e^t - e^{-t}$$
\end{solution}

\textbf{作业（5.3）}：p145：1（1、2）、2、4（1、2、3）
